% ============================================================
% 基礎演習Ⅱ　総まとめ冊子　共通プリアンブル
%   各冊子（summary-01-05.tex など）の冒頭で % ============================================================
% 基礎演習Ⅱ　総まとめ冊子　共通プリアンブル
%   各冊子（summary-01-05.tex など）の冒頭で % ============================================================
% 基礎演習Ⅱ　総まとめ冊子　共通プリアンブル
%   各冊子（summary-01-05.tex など）の冒頭で % ============================================================
% 基礎演習Ⅱ　総まとめ冊子　共通プリアンブル
%   各冊子（summary-01-05.tex など）の冒頭で \input{preamble} して使う。
%   LuaLaTeX（jlreq）で組む。色と部品はワークシート（../worksheets/preamble.tex）に合わせてある。
% ============================================================
\documentclass[paper=a4paper,fontsize=10pt,line_length=48zw,number_of_lines=44,
  baselineskip=1.6zw,head_space=20mm]{jlreq}

\usepackage{xcolor}
\usepackage[most]{tcolorbox}
\usepackage{tabularray}
\usepackage{enumitem}
\usepackage{needspace}
\usepackage{fancyhdr}
\usepackage[hidelinks]{hyperref}

% ------------------------------------------------------------
% 色（サイトの assets/css/style.css、ワークシートに合わせる）
% ------------------------------------------------------------
\definecolor{ink}{HTML}{1D2330}
\definecolor{inktwo}{HTML}{4A5261}
\definecolor{inkthree}{HTML}{757D8A}
\definecolor{rulegray}{HTML}{B9BEB4}
\definecolor{papertwo}{HTML}{ECEEE8}
\definecolor{pone}{HTML}{2747C9}  \definecolor{ponetint}{HTML}{DCE2F9}
\definecolor{ptwo}{HTML}{D68A17}  \definecolor{ptwotint}{HTML}{F6E4C3}
\definecolor{bad}{HTML}{B3261E}   \definecolor{badtint}{HTML}{F9DEDA}

% ------------------------------------------------------------
% 冊子ごとに設定する情報（各冊子の冒頭で \renewcommand する）
% ------------------------------------------------------------
\newcommand*\doctitle{総まとめ}
\newcommand*\docauthor{河合 勝彦}
\newcommand*\docaffil{名古屋市立大学 経済学部}
\newcommand*\docdate{}
\AtBeginDocument{\hypersetup{pdftitle={基礎演習Ⅱ　\doctitle},pdfauthor={\docauthor}}}

% ------------------------------------------------------------
% ページスタイル
% ------------------------------------------------------------
\pagestyle{fancy}
\fancyhf{}
\renewcommand{\headrulewidth}{0pt}
\fancyfoot[L]{\sffamily\scriptsize\color{inkthree}基礎演習Ⅱ　2026年度後期　\doctitle}
\fancyfoot[R]{\sffamily\scriptsize\color{inkthree}\thepage}
\fancypagestyle{cover}{\fancyhf{}}

\setlist{nosep,leftmargin=1.5em}
\setcounter{tocdepth}{2}
\widowpenalty=10000 \clubpenalty=10000

% ------------------------------------------------------------
% 見出し
% ------------------------------------------------------------
\ModifyHeading{section}{font=\large\sffamily\bfseries\color{pone},
  before_space=1.2\baselineskip,after_space=0.5\baselineskip,lines=3,
  label_format={\thesection},after_label_space=1\zw}
\ModifyHeading{subsection}{font=\normalsize\sffamily\bfseries\color{ink},
  before_space=0.9\baselineskip,after_space=0.3\baselineskip,lines=2,
  label_format={\thesubsection},after_label_space=0.8\zw}

% 部の帯　\partband{ラベル}{題}
\newcommand\partband[2]{%
  \clearpage
  \begin{tcolorbox}[enhanced,colback=pone,colframe=pone,arc=2mm,boxrule=0pt,
      left=4mm,right=4mm,top=2.5mm,bottom=2.5mm,fontupper=\color{white}\sffamily]
    {\small #1}\par\vspace{0.6mm}{\LARGE\bfseries #2}
  \end{tcolorbox}\vspace{1mm}}

% 付録に切り替える（節番号を 付録A、A.1 にする）
\newcommand\appendixmode{%
  \setcounter{section}{0}%
  \renewcommand\thesection{付録\Alph{section}}%
  \renewcommand\thesubsection{\Alph{section}.\arabic{subsection}}%
  \renewcommand\theHsection{appendix.\Alph{section}}%
  \renewcommand\theHsubsection{appendix.\Alph{section}.\arabic{subsection}}}

% 各回の見出し　\sessionhead{回}{日付}{題}{タグ}
\newcommand\sessionhead[4]{%
  \par\Needspace{34mm}\vspace{0.9\baselineskip}%
  \phantomsection\addcontentsline{toc}{subsection}{第#1回　#3}%
  \begin{tcolorbox}[enhanced,colback=ponetint,colframe=ponetint,arc=1.5mm,boxrule=0pt,
      borderline west={1.5mm}{0pt}{pone},left=4mm,right=3mm,top=1.5mm,bottom=1.5mm,
      before skip=0pt,after skip=2.5mm]
    {\sffamily\small\color{pone}第#1回　#2\if\relax\detokenize{#4}\relax\else　\textbar　#4\fi}\par
    \vspace{0.3mm}{\sffamily\large\bfseries\color{ink}#3}
  \end{tcolorbox}}

% ------------------------------------------------------------
% 囲み
% ------------------------------------------------------------
% 要点　\begin{keybox}{見出し} ... \end{keybox}
\newtcolorbox{keybox}[1]{enhanced,colback=papertwo,colframe=papertwo,boxrule=0pt,arc=1.5mm,
  left=3.5mm,right=3mm,top=1.5mm,bottom=1.5mm,fontupper=\small,
  before upper={{\sffamily\bfseries\color{pone}#1}\par\vspace{0.4mm}},
  before skip=2mm,after skip=2.5mm}

% 注意　\begin{caution}{見出し} ... \end{caution}
\newtcolorbox{caution}[1]{enhanced,colback=badtint,colframe=badtint,boxrule=0pt,arc=1.5mm,
  left=3.5mm,right=3mm,top=1.5mm,bottom=1.5mm,fontupper=\small,
  before upper={{\sffamily\bfseries\color{bad}#1}\par\vspace{0.4mm}},
  before skip=2mm,after skip=2.5mm}

% プロンプト　\begin{promptbox}[題] ... \end{promptbox}
%   （　）の部分は自分の内容に置き換える。行は空行で区切る。
\newtcolorbox{promptbox}[1][]{enhanced,colback=papertwo,colframe=papertwo,boxrule=0pt,
  arc=1mm,borderline west={1.2mm}{0pt}{pone},left=3.5mm,right=3mm,top=1.5mm,bottom=1.8mm,
  fontupper=\sffamily\small\color{ink}\parindent0pt,before skip=1.5mm,after skip=2mm,
  title={\sffamily\scriptsize\bfseries\color{pone}プロンプト\if\relax\detokenize{#1}\relax\else　#1\fi},
  attach title to upper={\par\vspace{0.6mm}}}

% チェック項目
\newenvironment{checks}{\begin{itemize}[label={□},leftmargin=1.5em,itemsep=0.8mm]\small}{\end{itemize}}

% ------------------------------------------------------------
% プロンプト集の部品
% ------------------------------------------------------------
% 項目の見出し　\promptitem{番号}{題}{使う場面}
\newcommand\promptitem[3]{%
  \par\Needspace{48mm}\vspace{1.2\baselineskip}%
  \phantomsection\label{p:#1}%
  \noindent{\sffamily\bfseries\colorbox{pone}{\strut\color{white}\,#1\,}}\enspace
  {\sffamily\bfseries\large #2}\par
  \vspace{0.4mm}\noindent{\small\color{inktwo}使う場面：#3}\par\vspace{0.6mm}}

% 注記　\pnote{ラベル}{本文}（ラベルは「解説」「続けて聞く」「確かめる」など）
\newcommand\pnote[2]{%
  \par\vspace{0.6mm}{\small\noindent\hangindent=5.5\zw\hangafter=1
  {\sffamily\bfseries\color{pone}\makebox[5.5\zw][l]{#1}}#2\par}}

% 授業で使ったプロンプトの印
\newcommand\usedmark{{\sffamily\scriptsize\color{white}\colorbox{ptwo}{\strut 授業で使用}}}

% ------------------------------------------------------------
% 表の共通スタイル
% ------------------------------------------------------------
\NewTblrEnviron{stab}
\SetTblrInner[stab]{hlines={0.4pt,rulegray},vlines={0.4pt,rulegray},
  cells={font=\small},rowsep=2.5pt,colsep=4pt,
  row{1}={bg=ponetint,font=\sffamily\footnotesize\bfseries}}
\newcommand\hc[1]{\SetCell{bg=ponetint,font=\sffamily\footnotesize\bfseries}#1}

% 付録の小見出し（次の項目と一緒に次の頁へ送る）
\newcommand\psubsection[1]{\par\Needspace{64mm}\subsection{#1}}
 して使う。
%   LuaLaTeX（jlreq）で組む。色と部品はワークシート（../worksheets/preamble.tex）に合わせてある。
% ============================================================
\documentclass[paper=a4paper,fontsize=10pt,line_length=48zw,number_of_lines=44,
  baselineskip=1.6zw,head_space=20mm]{jlreq}

\usepackage{xcolor}
\usepackage[most]{tcolorbox}
\usepackage{tabularray}
\usepackage{enumitem}
\usepackage{needspace}
\usepackage{fancyhdr}
\usepackage[hidelinks]{hyperref}

% ------------------------------------------------------------
% 色（サイトの assets/css/style.css、ワークシートに合わせる）
% ------------------------------------------------------------
\definecolor{ink}{HTML}{1D2330}
\definecolor{inktwo}{HTML}{4A5261}
\definecolor{inkthree}{HTML}{757D8A}
\definecolor{rulegray}{HTML}{B9BEB4}
\definecolor{papertwo}{HTML}{ECEEE8}
\definecolor{pone}{HTML}{2747C9}  \definecolor{ponetint}{HTML}{DCE2F9}
\definecolor{ptwo}{HTML}{D68A17}  \definecolor{ptwotint}{HTML}{F6E4C3}
\definecolor{bad}{HTML}{B3261E}   \definecolor{badtint}{HTML}{F9DEDA}

% ------------------------------------------------------------
% 冊子ごとに設定する情報（各冊子の冒頭で \renewcommand する）
% ------------------------------------------------------------
\newcommand*\doctitle{総まとめ}
\newcommand*\docauthor{河合 勝彦}
\newcommand*\docaffil{名古屋市立大学 経済学部}
\newcommand*\docdate{}
\AtBeginDocument{\hypersetup{pdftitle={基礎演習Ⅱ　\doctitle},pdfauthor={\docauthor}}}

% ------------------------------------------------------------
% ページスタイル
% ------------------------------------------------------------
\pagestyle{fancy}
\fancyhf{}
\renewcommand{\headrulewidth}{0pt}
\fancyfoot[L]{\sffamily\scriptsize\color{inkthree}基礎演習Ⅱ　2026年度後期　\doctitle}
\fancyfoot[R]{\sffamily\scriptsize\color{inkthree}\thepage}
\fancypagestyle{cover}{\fancyhf{}}

\setlist{nosep,leftmargin=1.5em}
\setcounter{tocdepth}{2}
\widowpenalty=10000 \clubpenalty=10000

% ------------------------------------------------------------
% 見出し
% ------------------------------------------------------------
\ModifyHeading{section}{font=\large\sffamily\bfseries\color{pone},
  before_space=1.2\baselineskip,after_space=0.5\baselineskip,lines=3,
  label_format={\thesection},after_label_space=1\zw}
\ModifyHeading{subsection}{font=\normalsize\sffamily\bfseries\color{ink},
  before_space=0.9\baselineskip,after_space=0.3\baselineskip,lines=2,
  label_format={\thesubsection},after_label_space=0.8\zw}

% 部の帯　\partband{ラベル}{題}
\newcommand\partband[2]{%
  \clearpage
  \begin{tcolorbox}[enhanced,colback=pone,colframe=pone,arc=2mm,boxrule=0pt,
      left=4mm,right=4mm,top=2.5mm,bottom=2.5mm,fontupper=\color{white}\sffamily]
    {\small #1}\par\vspace{0.6mm}{\LARGE\bfseries #2}
  \end{tcolorbox}\vspace{1mm}}

% 付録に切り替える（節番号を 付録A、A.1 にする）
\newcommand\appendixmode{%
  \setcounter{section}{0}%
  \renewcommand\thesection{付録\Alph{section}}%
  \renewcommand\thesubsection{\Alph{section}.\arabic{subsection}}%
  \renewcommand\theHsection{appendix.\Alph{section}}%
  \renewcommand\theHsubsection{appendix.\Alph{section}.\arabic{subsection}}}

% 各回の見出し　\sessionhead{回}{日付}{題}{タグ}
\newcommand\sessionhead[4]{%
  \par\Needspace{34mm}\vspace{0.9\baselineskip}%
  \phantomsection\addcontentsline{toc}{subsection}{第#1回　#3}%
  \begin{tcolorbox}[enhanced,colback=ponetint,colframe=ponetint,arc=1.5mm,boxrule=0pt,
      borderline west={1.5mm}{0pt}{pone},left=4mm,right=3mm,top=1.5mm,bottom=1.5mm,
      before skip=0pt,after skip=2.5mm]
    {\sffamily\small\color{pone}第#1回　#2\if\relax\detokenize{#4}\relax\else　\textbar　#4\fi}\par
    \vspace{0.3mm}{\sffamily\large\bfseries\color{ink}#3}
  \end{tcolorbox}}

% ------------------------------------------------------------
% 囲み
% ------------------------------------------------------------
% 要点　\begin{keybox}{見出し} ... \end{keybox}
\newtcolorbox{keybox}[1]{enhanced,colback=papertwo,colframe=papertwo,boxrule=0pt,arc=1.5mm,
  left=3.5mm,right=3mm,top=1.5mm,bottom=1.5mm,fontupper=\small,
  before upper={{\sffamily\bfseries\color{pone}#1}\par\vspace{0.4mm}},
  before skip=2mm,after skip=2.5mm}

% 注意　\begin{caution}{見出し} ... \end{caution}
\newtcolorbox{caution}[1]{enhanced,colback=badtint,colframe=badtint,boxrule=0pt,arc=1.5mm,
  left=3.5mm,right=3mm,top=1.5mm,bottom=1.5mm,fontupper=\small,
  before upper={{\sffamily\bfseries\color{bad}#1}\par\vspace{0.4mm}},
  before skip=2mm,after skip=2.5mm}

% プロンプト　\begin{promptbox}[題] ... \end{promptbox}
%   （　）の部分は自分の内容に置き換える。行は空行で区切る。
\newtcolorbox{promptbox}[1][]{enhanced,colback=papertwo,colframe=papertwo,boxrule=0pt,
  arc=1mm,borderline west={1.2mm}{0pt}{pone},left=3.5mm,right=3mm,top=1.5mm,bottom=1.8mm,
  fontupper=\sffamily\small\color{ink}\parindent0pt,before skip=1.5mm,after skip=2mm,
  title={\sffamily\scriptsize\bfseries\color{pone}プロンプト\if\relax\detokenize{#1}\relax\else　#1\fi},
  attach title to upper={\par\vspace{0.6mm}}}

% チェック項目
\newenvironment{checks}{\begin{itemize}[label={□},leftmargin=1.5em,itemsep=0.8mm]\small}{\end{itemize}}

% ------------------------------------------------------------
% プロンプト集の部品
% ------------------------------------------------------------
% 項目の見出し　\promptitem{番号}{題}{使う場面}
\newcommand\promptitem[3]{%
  \par\Needspace{48mm}\vspace{1.2\baselineskip}%
  \phantomsection\label{p:#1}%
  \noindent{\sffamily\bfseries\colorbox{pone}{\strut\color{white}\,#1\,}}\enspace
  {\sffamily\bfseries\large #2}\par
  \vspace{0.4mm}\noindent{\small\color{inktwo}使う場面：#3}\par\vspace{0.6mm}}

% 注記　\pnote{ラベル}{本文}（ラベルは「解説」「続けて聞く」「確かめる」など）
\newcommand\pnote[2]{%
  \par\vspace{0.6mm}{\small\noindent\hangindent=5.5\zw\hangafter=1
  {\sffamily\bfseries\color{pone}\makebox[5.5\zw][l]{#1}}#2\par}}

% 授業で使ったプロンプトの印
\newcommand\usedmark{{\sffamily\scriptsize\color{white}\colorbox{ptwo}{\strut 授業で使用}}}

% ------------------------------------------------------------
% 表の共通スタイル
% ------------------------------------------------------------
\NewTblrEnviron{stab}
\SetTblrInner[stab]{hlines={0.4pt,rulegray},vlines={0.4pt,rulegray},
  cells={font=\small},rowsep=2.5pt,colsep=4pt,
  row{1}={bg=ponetint,font=\sffamily\footnotesize\bfseries}}
\newcommand\hc[1]{\SetCell{bg=ponetint,font=\sffamily\footnotesize\bfseries}#1}

% 付録の小見出し（次の項目と一緒に次の頁へ送る）
\newcommand\psubsection[1]{\par\Needspace{64mm}\subsection{#1}}
 して使う。
%   LuaLaTeX（jlreq）で組む。色と部品はワークシート（../worksheets/preamble.tex）に合わせてある。
% ============================================================
\documentclass[paper=a4paper,fontsize=10pt,line_length=48zw,number_of_lines=44,
  baselineskip=1.6zw,head_space=20mm]{jlreq}

\usepackage{xcolor}
\usepackage[most]{tcolorbox}
\usepackage{tabularray}
\usepackage{enumitem}
\usepackage{needspace}
\usepackage{fancyhdr}
\usepackage[hidelinks]{hyperref}

% ------------------------------------------------------------
% 色（サイトの assets/css/style.css、ワークシートに合わせる）
% ------------------------------------------------------------
\definecolor{ink}{HTML}{1D2330}
\definecolor{inktwo}{HTML}{4A5261}
\definecolor{inkthree}{HTML}{757D8A}
\definecolor{rulegray}{HTML}{B9BEB4}
\definecolor{papertwo}{HTML}{ECEEE8}
\definecolor{pone}{HTML}{2747C9}  \definecolor{ponetint}{HTML}{DCE2F9}
\definecolor{ptwo}{HTML}{D68A17}  \definecolor{ptwotint}{HTML}{F6E4C3}
\definecolor{bad}{HTML}{B3261E}   \definecolor{badtint}{HTML}{F9DEDA}

% ------------------------------------------------------------
% 冊子ごとに設定する情報（各冊子の冒頭で \renewcommand する）
% ------------------------------------------------------------
\newcommand*\doctitle{総まとめ}
\newcommand*\docauthor{河合 勝彦}
\newcommand*\docaffil{名古屋市立大学 経済学部}
\newcommand*\docdate{}
\AtBeginDocument{\hypersetup{pdftitle={基礎演習Ⅱ　\doctitle},pdfauthor={\docauthor}}}

% ------------------------------------------------------------
% ページスタイル
% ------------------------------------------------------------
\pagestyle{fancy}
\fancyhf{}
\renewcommand{\headrulewidth}{0pt}
\fancyfoot[L]{\sffamily\scriptsize\color{inkthree}基礎演習Ⅱ　2026年度後期　\doctitle}
\fancyfoot[R]{\sffamily\scriptsize\color{inkthree}\thepage}
\fancypagestyle{cover}{\fancyhf{}}

\setlist{nosep,leftmargin=1.5em}
\setcounter{tocdepth}{2}
\widowpenalty=10000 \clubpenalty=10000

% ------------------------------------------------------------
% 見出し
% ------------------------------------------------------------
\ModifyHeading{section}{font=\large\sffamily\bfseries\color{pone},
  before_space=1.2\baselineskip,after_space=0.5\baselineskip,lines=3,
  label_format={\thesection},after_label_space=1\zw}
\ModifyHeading{subsection}{font=\normalsize\sffamily\bfseries\color{ink},
  before_space=0.9\baselineskip,after_space=0.3\baselineskip,lines=2,
  label_format={\thesubsection},after_label_space=0.8\zw}

% 部の帯　\partband{ラベル}{題}
\newcommand\partband[2]{%
  \clearpage
  \begin{tcolorbox}[enhanced,colback=pone,colframe=pone,arc=2mm,boxrule=0pt,
      left=4mm,right=4mm,top=2.5mm,bottom=2.5mm,fontupper=\color{white}\sffamily]
    {\small #1}\par\vspace{0.6mm}{\LARGE\bfseries #2}
  \end{tcolorbox}\vspace{1mm}}

% 付録に切り替える（節番号を 付録A、A.1 にする）
\newcommand\appendixmode{%
  \setcounter{section}{0}%
  \renewcommand\thesection{付録\Alph{section}}%
  \renewcommand\thesubsection{\Alph{section}.\arabic{subsection}}%
  \renewcommand\theHsection{appendix.\Alph{section}}%
  \renewcommand\theHsubsection{appendix.\Alph{section}.\arabic{subsection}}}

% 各回の見出し　\sessionhead{回}{日付}{題}{タグ}
\newcommand\sessionhead[4]{%
  \par\Needspace{34mm}\vspace{0.9\baselineskip}%
  \phantomsection\addcontentsline{toc}{subsection}{第#1回　#3}%
  \begin{tcolorbox}[enhanced,colback=ponetint,colframe=ponetint,arc=1.5mm,boxrule=0pt,
      borderline west={1.5mm}{0pt}{pone},left=4mm,right=3mm,top=1.5mm,bottom=1.5mm,
      before skip=0pt,after skip=2.5mm]
    {\sffamily\small\color{pone}第#1回　#2\if\relax\detokenize{#4}\relax\else　\textbar　#4\fi}\par
    \vspace{0.3mm}{\sffamily\large\bfseries\color{ink}#3}
  \end{tcolorbox}}

% ------------------------------------------------------------
% 囲み
% ------------------------------------------------------------
% 要点　\begin{keybox}{見出し} ... \end{keybox}
\newtcolorbox{keybox}[1]{enhanced,colback=papertwo,colframe=papertwo,boxrule=0pt,arc=1.5mm,
  left=3.5mm,right=3mm,top=1.5mm,bottom=1.5mm,fontupper=\small,
  before upper={{\sffamily\bfseries\color{pone}#1}\par\vspace{0.4mm}},
  before skip=2mm,after skip=2.5mm}

% 注意　\begin{caution}{見出し} ... \end{caution}
\newtcolorbox{caution}[1]{enhanced,colback=badtint,colframe=badtint,boxrule=0pt,arc=1.5mm,
  left=3.5mm,right=3mm,top=1.5mm,bottom=1.5mm,fontupper=\small,
  before upper={{\sffamily\bfseries\color{bad}#1}\par\vspace{0.4mm}},
  before skip=2mm,after skip=2.5mm}

% プロンプト　\begin{promptbox}[題] ... \end{promptbox}
%   （　）の部分は自分の内容に置き換える。行は空行で区切る。
\newtcolorbox{promptbox}[1][]{enhanced,colback=papertwo,colframe=papertwo,boxrule=0pt,
  arc=1mm,borderline west={1.2mm}{0pt}{pone},left=3.5mm,right=3mm,top=1.5mm,bottom=1.8mm,
  fontupper=\sffamily\small\color{ink}\parindent0pt,before skip=1.5mm,after skip=2mm,
  title={\sffamily\scriptsize\bfseries\color{pone}プロンプト\if\relax\detokenize{#1}\relax\else　#1\fi},
  attach title to upper={\par\vspace{0.6mm}}}

% チェック項目
\newenvironment{checks}{\begin{itemize}[label={□},leftmargin=1.5em,itemsep=0.8mm]\small}{\end{itemize}}

% ------------------------------------------------------------
% プロンプト集の部品
% ------------------------------------------------------------
% 項目の見出し　\promptitem{番号}{題}{使う場面}
\newcommand\promptitem[3]{%
  \par\Needspace{48mm}\vspace{1.2\baselineskip}%
  \phantomsection\label{p:#1}%
  \noindent{\sffamily\bfseries\colorbox{pone}{\strut\color{white}\,#1\,}}\enspace
  {\sffamily\bfseries\large #2}\par
  \vspace{0.4mm}\noindent{\small\color{inktwo}使う場面：#3}\par\vspace{0.6mm}}

% 注記　\pnote{ラベル}{本文}（ラベルは「解説」「続けて聞く」「確かめる」など）
\newcommand\pnote[2]{%
  \par\vspace{0.6mm}{\small\noindent\hangindent=5.5\zw\hangafter=1
  {\sffamily\bfseries\color{pone}\makebox[5.5\zw][l]{#1}}#2\par}}

% 授業で使ったプロンプトの印
\newcommand\usedmark{{\sffamily\scriptsize\color{white}\colorbox{ptwo}{\strut 授業で使用}}}

% ------------------------------------------------------------
% 表の共通スタイル
% ------------------------------------------------------------
\NewTblrEnviron{stab}
\SetTblrInner[stab]{hlines={0.4pt,rulegray},vlines={0.4pt,rulegray},
  cells={font=\small},rowsep=2.5pt,colsep=4pt,
  row{1}={bg=ponetint,font=\sffamily\footnotesize\bfseries}}
\newcommand\hc[1]{\SetCell{bg=ponetint,font=\sffamily\footnotesize\bfseries}#1}

% 付録の小見出し（次の項目と一緒に次の頁へ送る）
\newcommand\psubsection[1]{\par\Needspace{64mm}\subsection{#1}}
 して使う。
%   LuaLaTeX（jlreq）で組む。色と部品はワークシート（../worksheets/preamble.tex）に合わせてある。
% ============================================================
\documentclass[paper=a4paper,fontsize=10pt,line_length=48zw,number_of_lines=44,
  baselineskip=1.6zw,head_space=20mm]{jlreq}

\usepackage{xcolor}
\usepackage[most]{tcolorbox}
\usepackage{tabularray}
\usepackage{enumitem}
\usepackage{needspace}
\usepackage{fancyhdr}
\usepackage[hidelinks]{hyperref}

% ------------------------------------------------------------
% 色（サイトの assets/css/style.css、ワークシートに合わせる）
% ------------------------------------------------------------
\definecolor{ink}{HTML}{1D2330}
\definecolor{inktwo}{HTML}{4A5261}
\definecolor{inkthree}{HTML}{757D8A}
\definecolor{rulegray}{HTML}{B9BEB4}
\definecolor{papertwo}{HTML}{ECEEE8}
\definecolor{pone}{HTML}{2747C9}  \definecolor{ponetint}{HTML}{DCE2F9}
\definecolor{ptwo}{HTML}{D68A17}  \definecolor{ptwotint}{HTML}{F6E4C3}
\definecolor{bad}{HTML}{B3261E}   \definecolor{badtint}{HTML}{F9DEDA}

% ------------------------------------------------------------
% 冊子ごとに設定する情報（各冊子の冒頭で \renewcommand する）
% ------------------------------------------------------------
\newcommand*\doctitle{総まとめ}
\newcommand*\docauthor{河合 勝彦}
\newcommand*\docaffil{名古屋市立大学 経済学部}
\newcommand*\docdate{}
\AtBeginDocument{\hypersetup{pdftitle={基礎演習Ⅱ　\doctitle},pdfauthor={\docauthor}}}

% ------------------------------------------------------------
% ページスタイル
% ------------------------------------------------------------
\pagestyle{fancy}
\fancyhf{}
\renewcommand{\headrulewidth}{0pt}
\fancyfoot[L]{\sffamily\scriptsize\color{inkthree}基礎演習Ⅱ　2026年度後期　\doctitle}
\fancyfoot[R]{\sffamily\scriptsize\color{inkthree}\thepage}
\fancypagestyle{cover}{\fancyhf{}}

\setlist{nosep,leftmargin=1.5em}
\setcounter{tocdepth}{2}
\widowpenalty=10000 \clubpenalty=10000

% ------------------------------------------------------------
% 見出し
% ------------------------------------------------------------
\ModifyHeading{section}{font=\large\sffamily\bfseries\color{pone},
  before_space=1.2\baselineskip,after_space=0.5\baselineskip,lines=3,
  label_format={\thesection},after_label_space=1\zw}
\ModifyHeading{subsection}{font=\normalsize\sffamily\bfseries\color{ink},
  before_space=0.9\baselineskip,after_space=0.3\baselineskip,lines=2,
  label_format={\thesubsection},after_label_space=0.8\zw}

% 部の帯　\partband{ラベル}{題}
\newcommand\partband[2]{%
  \clearpage
  \begin{tcolorbox}[enhanced,colback=pone,colframe=pone,arc=2mm,boxrule=0pt,
      left=4mm,right=4mm,top=2.5mm,bottom=2.5mm,fontupper=\color{white}\sffamily]
    {\small #1}\par\vspace{0.6mm}{\LARGE\bfseries #2}
  \end{tcolorbox}\vspace{1mm}}

% 付録に切り替える（節番号を 付録A、A.1 にする）
\newcommand\appendixmode{%
  \setcounter{section}{0}%
  \renewcommand\thesection{付録\Alph{section}}%
  \renewcommand\thesubsection{\Alph{section}.\arabic{subsection}}%
  \renewcommand\theHsection{appendix.\Alph{section}}%
  \renewcommand\theHsubsection{appendix.\Alph{section}.\arabic{subsection}}}

% 各回の見出し　\sessionhead{回}{日付}{題}{タグ}
\newcommand\sessionhead[4]{%
  \par\Needspace{34mm}\vspace{0.9\baselineskip}%
  \phantomsection\addcontentsline{toc}{subsection}{第#1回　#3}%
  \begin{tcolorbox}[enhanced,colback=ponetint,colframe=ponetint,arc=1.5mm,boxrule=0pt,
      borderline west={1.5mm}{0pt}{pone},left=4mm,right=3mm,top=1.5mm,bottom=1.5mm,
      before skip=0pt,after skip=2.5mm]
    {\sffamily\small\color{pone}第#1回　#2\if\relax\detokenize{#4}\relax\else　\textbar　#4\fi}\par
    \vspace{0.3mm}{\sffamily\large\bfseries\color{ink}#3}
  \end{tcolorbox}}

% ------------------------------------------------------------
% 囲み
% ------------------------------------------------------------
% 要点　\begin{keybox}{見出し} ... \end{keybox}
\newtcolorbox{keybox}[1]{enhanced,colback=papertwo,colframe=papertwo,boxrule=0pt,arc=1.5mm,
  left=3.5mm,right=3mm,top=1.5mm,bottom=1.5mm,fontupper=\small,
  before upper={{\sffamily\bfseries\color{pone}#1}\par\vspace{0.4mm}},
  before skip=2mm,after skip=2.5mm}

% 注意　\begin{caution}{見出し} ... \end{caution}
\newtcolorbox{caution}[1]{enhanced,colback=badtint,colframe=badtint,boxrule=0pt,arc=1.5mm,
  left=3.5mm,right=3mm,top=1.5mm,bottom=1.5mm,fontupper=\small,
  before upper={{\sffamily\bfseries\color{bad}#1}\par\vspace{0.4mm}},
  before skip=2mm,after skip=2.5mm}

% プロンプト　\begin{promptbox}[題] ... \end{promptbox}
%   （　）の部分は自分の内容に置き換える。行は空行で区切る。
\newtcolorbox{promptbox}[1][]{enhanced,colback=papertwo,colframe=papertwo,boxrule=0pt,
  arc=1mm,borderline west={1.2mm}{0pt}{pone},left=3.5mm,right=3mm,top=1.5mm,bottom=1.8mm,
  fontupper=\sffamily\small\color{ink}\parindent0pt,before skip=1.5mm,after skip=2mm,
  title={\sffamily\scriptsize\bfseries\color{pone}プロンプト\if\relax\detokenize{#1}\relax\else　#1\fi},
  attach title to upper={\par\vspace{0.6mm}}}

% チェック項目
\newenvironment{checks}{\begin{itemize}[label={□},leftmargin=1.5em,itemsep=0.8mm]\small}{\end{itemize}}

% ------------------------------------------------------------
% プロンプト集の部品
% ------------------------------------------------------------
% 項目の見出し　\promptitem{番号}{題}{使う場面}
\newcommand\promptitem[3]{%
  \par\Needspace{48mm}\vspace{1.2\baselineskip}%
  \phantomsection\label{p:#1}%
  \noindent{\sffamily\bfseries\colorbox{pone}{\strut\color{white}\,#1\,}}\enspace
  {\sffamily\bfseries\large #2}\par
  \vspace{0.4mm}\noindent{\small\color{inktwo}使う場面：#3}\par\vspace{0.6mm}}

% 注記　\pnote{ラベル}{本文}（ラベルは「解説」「続けて聞く」「確かめる」など）
\newcommand\pnote[2]{%
  \par\vspace{0.6mm}{\small\noindent\hangindent=5.5\zw\hangafter=1
  {\sffamily\bfseries\color{pone}\makebox[5.5\zw][l]{#1}}#2\par}}

% 授業で使ったプロンプトの印
\newcommand\usedmark{{\sffamily\scriptsize\color{white}\colorbox{ptwo}{\strut 授業で使用}}}

% ------------------------------------------------------------
% 表の共通スタイル
% ------------------------------------------------------------
\NewTblrEnviron{stab}
\SetTblrInner[stab]{hlines={0.4pt,rulegray},vlines={0.4pt,rulegray},
  cells={font=\small},rowsep=2.5pt,colsep=4pt,
  row{1}={bg=ponetint,font=\sffamily\footnotesize\bfseries}}
\newcommand\hc[1]{\SetCell{bg=ponetint,font=\sffamily\footnotesize\bfseries}#1}

% 付録の小見出し（次の項目と一緒に次の頁へ送る）
\newcommand\psubsection[1]{\par\Needspace{64mm}\subsection{#1}}
